%=============================================================================!
% 2.3  MASS AND THE SECOND LAW                                                !
%=============================================================================!
\section{Mass and the second law}
\label{sec:ne-second}

The first law says what happens when no net force acts. The second says
what a net force does, and to state it we need one more quantity, mass,
which measures how strongly a body resists being accelerated.

\subsection*{Mass}

Pull two different bodies on a smooth level floor with the same force, for instance by stretching
the same spring by the same amount, and they accelerate differently: a
loaded trolley gathers speed more slowly than an empty one. Experiment
shows that the ratio of their accelerations, $a_2/a_1$, does not depend on
which force is used. Because the ratio does not depend on the force, it
belongs to the two bodies alone, and we can use it to give each body a
number, its mass, by defining the ratio of their masses as
\begin{equation*}
   \frac{m_1}{m_2} = \frac{a_2}{a_1} .
\end{equation*}
The body that accelerates less has the larger mass. If, under the same
pull, an empty trolley accelerates at \SI{3}{\metre\per\second\squared}
and a loaded one at \SI{1}{\metre\per\second\squared}, the loaded trolley
has three times the mass of the empty one. Choosing one body as a
standard, the kilogram (\si{\kilogram}), then fixes every other mass. Mass
is a new kind of quantity alongside length and time, and we give it the
dimension $\mathsf{M}$. It measures the inertia of the body, and
experiment also shows that when two bodies are joined their masses add.

\subsection*{The second law}

Two experiments now combine into a law. The mass
experiment says that one force gives two bodies accelerations in the
inverse ratio of their masses, $a_2/a_1 = m_1/m_2$, so the product $m_1
a_1 = m_2 a_2$ is the same for both bodies: it depends on the force and
not on the body it acts on. Experiment also shows that $n$ standard
springs pulling together give a body $n$ times the acceleration that one
gives it, so the product is $n$ times as large when the force is $n$
units. The product of mass and acceleration is therefore proportional to
the force, $F = c\,ma$
for some constant $c$, and the acceleration points along the force. We
are free to choose the unit of force, and choosing it to make $c = 1$, as
the newton below does, gives Newton's second law,
\begin{equation}
   \vb{F} = m\vb{a} ,
   \label{eq:ne-second-law}
\end{equation}
where $\vb{F}$ is the net force on the body. Because it is a vector
equation, it holds for each component separately: $F_x = m a_x$, $F_y =
m a_y$ and $F_z = m a_z$.

The law also fixes the unit of force. One newton (\si{\newton}) is the
force that gives one kilogram an acceleration of one metre per second per
second, $\SI{1}{\newton} = \SI{1}{\kilogram\metre\per\second\squared}$,
so force has dimension $\mathsf{M}\mathsf{L}\mathsf{T}^{-2}$. The spring
unit of \cref{sec:ne-force} was a stopgap: from now on forces are measured
in newtons, and spring balances are marked in newtons. A trolley of
\SI{20}{\kilogram} pushed with a net force of \SI{10}{\newton}
accelerates at $10/20 = \SI{0.5}{\metre\per\second\squared}$, in the
direction of the push.

\begin{keyidea}
In an inertial frame of reference the net force on a body equals its mass
times its acceleration, $\vb{F} = m\vb{a}$, where the net force is the
vector sum of every force acting on the body.
\end{keyidea}

\begin{selfcheck}
A net force of \SI{6}{\newton} acts on a \SI{2}{\kilogram} body. What is
its acceleration? What would it be if the mass were doubled?
\end{selfcheck}
